% ------------------------------------------------------------------------
%
\begin{frame}{Analysis of algorithms}

  The mathematical study of the efficiency of algorithms has been
  pioneered by Knuth, who called it \textbf{analysis of algorithms}.

  \bigskip

  Given a function definition,
  \begin{enumerate}

    \item we define a measure on the arguments, which represents their
      size;

    \item we define a measure on time, which abstracts the wall clock
      time;

    \item we express the abstract time needed to compute calls to that
      function in terms of the size of its arguments: this is the
      \textbf{cost function}.

  \end{enumerate}

  The lower the cost, the higher the efficiency.

  For example, for sorting objects, called \textbf{keys}, by comparing
  them,
  \begin{enumerate}

    \item the input size is the number of keys,

    \item the abstract unit of time is one comparison,

    \item and the cost maps the number of keys to the number of
      comparisons to sort the keys.

  \end{enumerate}

\end{frame}

% ------------------------------------------------------------------------
%
\begin{frame}{Exact cost}

  Rewrite systems enable a rather natural notion of cost for
  functional programs: it is the number of rewrites to reach the value
  of a function call, assuming that the arguments are values.

  \bigskip

  In other words, the \textbf{exact cost} is the number of calls
  needed to compute the value of a given call.

  \bigskip

  For stacks, the measure of the size of the input is the number of
  items it contains, or \textbf{length} of the stack.

\end{frame}

% ------------------------------------------------------------------------
%
\begin{frame}{Exact cost / Appending a stack}

  For instance, let us recall the appending of a stack:
  \begin{align*}
    \fun{cat}(        \el,t) &\xrightarrow{\alpha} t;\\
    \fun{cat}(\cons{x}{s},t) &\xrightarrow{\smash{\beta}}
    \cons{x}{\fun{cat}(s,t)}.
  \end{align*}

  We observe that~\(t\) is invariant, so the cost depends only on the
  size of the first argument.

  \bigskip

  Let \(\C{\fun{cat}}{n}\) be the cost of the call \(\fun{cat}(s,t)\),
  for any~\(t\), where \(n\)~is the size of~\(s\). Each rule yields an
  equation on their cost:
  \begin{align*}
    \C{\fun{cat}}{0} & \eqn{\alpha} 1,\\
    \C{\fun{cat}}{n+1} & \eqn{\smash{\beta}} 1 + \C{\fun{cat}}{n}.
  \end{align*}

  Together, they yield \(\C{\fun{cat}}{n} = n + 1.\) Why?

\end{frame}

% ------------------------------------------------------------------------
%
\begin{frame}{Why?}

  The trick consists in making the difference of two successive terms
  in the series and then summing up these differences:
  \begin{align*}
    \C{\fun{cat}}{n+1} - \C{\fun{cat}}{n} &= 1 & \text{where} \;\, n
    \geqslant 0\\
    \sum_{n=0}^{p}{(\C{\fun{cat}}{n+1} - \C{\fun{cat}}{n})} &=
    \sum_{n=0}^{p}{1} & \text{where} \;\, p \geqslant 0\\
    \C{\fun{cat}}{p+1} - \C{\fun{cat}}{0} &= p + 1\\
    \intertext{Replacing \(p+1\) by \(n\) we finally get}
    \C{\fun{cat}}{n} - \C{\fun{cat}}{0} &= n & \text{where}\;\, n > 0\\
    \C{\fun{cat}}{n}  &= n + 1
  \end{align*}
  Since that replacing \(n\) by \(0\) in the latter formula gives
  \(\C{\fun{cat}}{0} = 0 + 1 = 1\), we can have only one formula for
  the cost: \(\C{\fun{cat}}{n} = n + 1,\) where \(n \geqslant
  0.\)\hfill\(\Box\)

\end{frame}

% ------------------------------------------------------------------------
%
\begin{frame}{Exact cost / Reversal of a list}

  Let us consider the definition of a function \fun{rev\(_0\)}
  reversing a stack:
  \begin{equation*}
    \begin{array}{@{}r@{\;}l@{\;}lr@{\;}l@{\;}l@{}}
      \fun{cat}(\el,t)
      & \xrightarrow{\alpha} & t;
      & \fun{rev}_0\,\el
      & \xrightarrow{\smash{\gamma}} & \el;\\
      \fun{cat}(\cons{x}{s},t)
      & \xrightarrow{\smash{\beta}} & \cons{x}{\fun{cat}(s,t)}.
      & \fun{rev}_0(\cons{x}{s})
      & \xrightarrow{\smash{\delta}} & \fun{cat}(\fun{rev}_0\,s, [x]).
    \end{array}
  \end{equation*}

  The evaluation of \(\fun{rev}_0\,[3;2;1]\) is shown with abstract
  syntax trees:
  \begin{center}
    \includegraphics[bb=71 637 350 720]{rev0_321_0}
  \end{center}

\end{frame}

% ------------------------------------------------------------------------
%
\begin{frame}{Exact cost / Reversal of a list}

  Here are the definitions again:
  \begin{equation*}
    \begin{array}{@{}r@{\;}l@{\;}lr@{\;}l@{\;}l@{}}
      \fun{cat}(\el,t)
      & \xrightarrow{\alpha} & t;
      & \fun{rev}_0\,\el
      & \xrightarrow{\smash{\gamma}} & \el;\\
      \fun{cat}(\cons{x}{s},t)
      & \xrightarrow{\smash{\beta}} & \cons{x}{\fun{cat}(s,t)}.
      & \fun{rev}_0(\cons{x}{s})
      & \xrightarrow{\smash{\delta}} & \fun{cat}(\fun{rev}_0\,s, [x]).
    \end{array}
  \end{equation*}

  Let us resume the evaluation of \(\fun{rev}_0\,[3;2;1]\) from where
  we left:
  \begin{center}
    \includegraphics[bb=71 637 358 721]{rev0_321_1}
  \end{center}

\end{frame}

% ------------------------------------------------------------------------
%
\begin{frame}{Exact cost / Reversal of a list}

  Here are the definitions again:
  \begin{equation*}
    \begin{array}{@{}r@{\;}l@{\;}lr@{\;}l@{\;}l@{}}
      \fun{cat}(\el,t)
      & \xrightarrow{\alpha} & t;
      & \fun{rev}_0\,\el
      & \xrightarrow{\smash{\gamma}} & \el;\\
      \fun{cat}(\cons{x}{s},t)
      & \xrightarrow{\smash{\beta}} & \cons{x}{\fun{cat}(s,t)}.
      & \fun{rev}_0(\cons{x}{s})
      & \xrightarrow{\smash{\delta}} & \fun{cat}(\fun{rev}_0\,s, [x]).
    \end{array}
  \end{equation*}

  Let us resume the evaluation of \(\fun{rev}_0\,[3;2;1]\) from where
  we left:
  \begin{center}
    \includegraphics[bb=71 639 360 721]{rev0_321_1_1}
  \end{center}

\end{frame}

% ------------------------------------------------------------------------
%
\begin{frame}{Exact cost / Reversal of a list}

  Here are the definitions again:
  \begin{equation*}
    \begin{array}{@{}r@{\;}l@{\;}lr@{\;}l@{\;}l@{}}
      \fun{cat}(\el,t)
      & \xrightarrow{\alpha} & t;
      & \fun{rev}_0\,\el
      & \xrightarrow{\smash{\gamma}} & \el;\\
      \fun{cat}(\cons{x}{s},t)
      & \xrightarrow{\smash{\beta}} & \cons{x}{\fun{cat}(s,t)}.
      & \fun{rev}_0(\cons{x}{s})
      & \xrightarrow{\smash{\delta}} & \fun{cat}(\fun{rev}_0\,s, [x]).
    \end{array}
  \end{equation*}

  Let us resume the evaluation of \(\fun{rev}_0\,[3;2;1]\) from where
  we left:
  \begin{center}
    \includegraphics[bb=71 639 334 721]{rev0_321_2}
  \end{center}

\end{frame}

% ------------------------------------------------------------------------
%
\begin{frame}{Exact cost / Reversal of a list}

  Here are the definitions again:
  \begin{equation*}
    \begin{array}{@{}r@{\;}l@{\;}lr@{\;}l@{\;}l@{}}
      \fun{cat}(\el,t)
      & \xrightarrow{\alpha} & t;
      & \fun{rev}_0\,\el
      & \xrightarrow{\smash{\gamma}} & \el;\\
      \fun{cat}(\cons{x}{s},t)
      & \xrightarrow{\smash{\beta}} & \cons{x}{\fun{cat}(s,t)}.
      & \fun{rev}_0(\cons{x}{s})
      & \xrightarrow{\smash{\delta}} & \fun{cat}(\fun{rev}_0\,s, [x]).
    \end{array}
  \end{equation*}

  Let us resume the evaluation of \(\fun{rev}_0\,[3;2;1]\) from where
  we left:
  \begin{center}
    \includegraphics[bb=71 637 372 721]{rev0_321_2_1}
  \end{center}

\end{frame}

% ------------------------------------------------------------------------
%
\begin{frame}{Exact cost / Reversal of a list}

  Here are the definitions again:
  \begin{equation*}
    \begin{array}{@{}r@{\;}l@{\;}lr@{\;}l@{\;}l@{}}
      \fun{cat}(\el,t)
      & \xrightarrow{\alpha} & t;
      & \fun{rev}_0\,\el
      & \xrightarrow{\smash{\gamma}} & \el;\\
      \fun{cat}(\cons{x}{s},t)
      & \xrightarrow{\smash{\beta}} & \cons{x}{\fun{cat}(s,t)}.
      & \fun{rev}_0(\cons{x}{s})
      & \xrightarrow{\smash{\delta}} & \fun{cat}(\fun{rev}_0\,s, [x]).
    \end{array}
  \end{equation*}

  Let us resume the evaluation of \(\fun{rev}_0\,[3;2;1]\) from where
  we left:
  \begin{center}
    \includegraphics[bb=71 637 390 721]{rev0_321_3}
  \end{center}

\end{frame}

% ------------------------------------------------------------------------
%
\begin{frame}{Exact cost / Reversal of a list}

  \begin{itemize}

    \item Here are the definitions again:
      \begin{equation*}
        \begin{array}{@{}r@{\;}l@{\;}lr@{\;}l@{\;}l@{}}
          \fun{cat}(\el,t)
          & \xrightarrow{\alpha} & t;
          & \fun{rev}_0\,\el
          & \xrightarrow{\smash{\gamma}} & \el;\\
          \fun{cat}(\cons{x}{s},t)
          & \xrightarrow{\smash{\beta}} & \cons{x}{\fun{cat}(s,t)}.
          & \fun{rev}_0(\cons{x}{s})
          & \xrightarrow{\smash{\delta}} & \fun{cat}(\fun{rev}_0\,s, [x]).
        \end{array}
      \end{equation*}

    \item Let \(\C{\fun{rev}_0}{k}\) be the cost of evaluating a call
      \(\fun{rev}_0(l)\), where the list~\(l\) has length~\(k\).

    \item Let \(\C{\fun{cat}}{k}\) be the cost of evaluating a call
      \(\fun{cat}(l)\), where the list~\(l\) has length~\(k\).

    \item The definition of \fun{rev\(_0\)} directly leads to the
      recurrence equations
      \begin{align*}
        \C{\fun{rev}_0}{0}   & \eqn{\gamma} 1,\\
        \C{\fun{rev}_0}{k+1} & \eqn{\smash{\delta}} 1 +
        \C{\fun{rev}_0}{k} + \C{\fun{cat}}{k} = \C{\fun{rev}_0}{k} + k
        + 2.
      \end{align*}
      because the length of \(\fun{rev}_0(s)\) is~\(k\) if the length
      of~\(s\) is~\(k\), and we already know \(\C{\fun{cat}}{k} = k +
      1\).

  \end{itemize}

\end{frame}

% ------------------------------------------------------------------------
%
\begin{frame}{Exact cost / Reversal of a list}

  We have
  \begin{equation*}
    \sum_{k=0}^{n-1}(\C{\fun{rev}_0}{k+1} - \C{\fun{rev}_0}{k})
    = \C{\fun{rev}_0}{n} - \C{\fun{rev}_0}{0}
    = \sum_{k=0}^{n-1}(k+2)
    = 2n + \sum_{k=0}^{n-1}{k}.
  \end{equation*}

  The remaining sum is a classic of algebra:
  \begin{equation*}
    2 \cdot \sum_{k=0}^{n-1}{k}
    = \sum_{k=0}^{n-1}{k} + \sum_{k=0}^{n-1}{k}
    = \sum_{k=0}^{n-1}{k} + \sum_{k=0}^{n-1}(n-k-1)
    = n(n-1).
  \end{equation*}

  Consequently,
  \begin{equation*}
    \sum_{k=0}^{n-1}{k} = \frac{1}{2} n(n-1).
  \end{equation*}

  We can finally conclude
  \begin{equation*}
    \C{\fun{rev}_0}{n}
    = \C{\fun{rev}_0}{0} + 2n + \frac{1}{2} n(n-1)
    = \frac{1}{2}n^2 + \frac{3}{2}n + 1
    \sim \frac{1}{2}n^2.
  \end{equation*}

\end{frame}

% ------------------------------------------------------------------------
%
\begin{frame}{Exact cost / Evaluation traces}

  Another way to reach the result is to induce an \textbf{evaluation
    trace}. A trace is a composition of rewrite rules, noted using the
  mathematical convention for multiplication.

  \bigskip

  From the evaluation of \(\fun{rev}_0\,[3;3;1]\), we infer the trace
  \(\T{\fun{rev}_0}{n}\) of the evaluation of \(\fun{rev}_0(s)\),
  where \(n\)~is the length of~\(s\):
  \begin{equation*}
    \T{\fun{rev}_0}{n} \; \text{:=} \;
    \delta^n\gamma\alpha(\beta\alpha)\dots(\beta^{n-1}\alpha)
    = \delta^n\gamma \prod_{k=0}^{n-1}{\beta^k\alpha}.
  \end{equation*}

  This reads: `Apply \(\delta\) \(n\)~times, then \(\gamma\), then
  \(\alpha\) etc.'

  \bigskip

  If we note \(\len{\T{\fun{rev}_0}{n}}\) the length
  of~\(\T{\fun{rev}_0}{n}\), that is, the number of rule applications
  it contains, we expect to have the equations \(\len{x} = 1\), for a
  rule~\(x\), and \(\len{x \cdot y} = \len{x} + \len{y}\), for rules
  \(x\)~and~\(y\).

\end{frame}

% ------------------------------------------------------------------------
%
\begin{frame}{Exact cost / Evaluation traces}

  By definition of the cost:
  \begin{align*}
    \C{\fun{rev}_0}{n}
    & \; \text{:=} \; \len{\T{\fun{rev}_0}{n}}
    = \left\lvert\delta^n\gamma \prod_{k=0}^{n-1}{\beta^k\alpha}\right\lvert
    = \len{\delta^n\gamma} + \sum_{k=0}^{n-1}\len{\beta^k\alpha}\\
    & = (n+1) + \!\sum_{k=0}^{n-1}(k+1)
      = (n+1) + \!\sum_{k=1}^{n}k
      =  \frac{1}{2}n^2 + \frac{3}{2}n + 1
      \sim  \frac{1}{2}n^2.
  \end{align*}
  Since we inferred our trace from an example, the cost function we
  derived may be wrong: we need to prove it by \textbf{induction}.

  \medskip

  Let \(\pred{Quad}{n}\) be the property \(\C{\fun{rev}_0}{n} = (n^2 +
  3n + 2)/2\). We already know that \(\pred{Quad}{0}\) is true. Let us
  suppose \(\pred{Quad}{n}\) is true for some value of~\(n\)
  (induction hypothesis) and let us prove \(\pred{Quad}{n+1}\).

  \medskip

  We know \(\C{\fun{rev}_0}{n+1} = \C{\fun{rev}_0}{n} + n + 2\). The
  induction hypothesis implies
  \begin{equation*}
    \C{\fun{rev}_0}{n+1}
    = (n^2 + 3n + 2)/2 + n + 2
    = ((n+1)^2 + 3(n+1) + 2)/2,
  \end{equation*}
  which is \(\pred{Quad}{n+1}\). Therefore, the induction principle
  says that the cost we found is always correct.\hfill\(\Box\)

\end{frame}

% ------------------------------------------------------------------------
%
\begin{frame}{Exact cost / Empirical approach}

  Yet another way to determine the cost of \fun{rev\(_0\)} consists in
  first \textbf{guessing} that it is quadratic, that is,
  \(\C{\fun{rev}_0}{n} = an^2 + bn + c\), where \(a\), \(b\) and~\(c\)
  are unknowns.

  \bigskip

  Since there are three coefficients, we only need three values
  of~\(\C{\fun{rev}_0}{n}\) to determine them, for instance \(n = 0,
  1, 2\).

  \bigskip

  Making some traces, we find \(\C{\fun{rev}_0}{0} = 1\),
  \(\C{\fun{rev}_0}{1} = 3\) and \(\C{\fun{rev}_0}{2} = 6\), so we
  solve the linear system
  \begin{equation*}
    \C{\fun{rev}_0}{0} = c = 1,\;
    \C{\fun{rev}_0}{1} = a + b + c = 3,\;
    \C{\fun{rev}_0}{2} = a \cdot 2^2 + b \cdot 2 + c = 6.
  \end{equation*}

  Hence \(a = \myfrac1/2\), \(b = \myfrac3/2\) and~\(c = 1\), that is,
  \(\C{\fun{rev}_0}{n} = (n^2 + 3n + 2)/2\).

  We may have been wrong with our guess of a quadratic cost: we need
  also here to prove it by \textbf{induction} exactly as before.

\end{frame}

% ------------------------------------------------------------------------
%
\begin{frame}{Exact cost / Efficient reversal}

  \label{def:rev}

  We would be remiss not to consider a more efficient (that is, less
  costly) definition of a reversal function. Here it is:
  \begin{align*}
    \fun{rev}\,s &\smashedrightarrow{\epsilon} \fun{rcat}(s,\el).\\
    \fun{rcat}(\el,t) & \smashedrightarrow{\zeta} t;\\
    \fun{rcat}(\cons{x}{s},t) &\smashedrightarrow{\eta}
    \fun{rcat}(s,\cons{x}{t}).
  \end{align*}
  (The name \fun{rcat} stands for \emph{reverse and catenate}.)

  \bigskip

  The additional parameter introduced by \fun{rcat} accumulates
  partial results, thus is called an \textbf{accumulator}:
  \begin{equation*}
    \begin{array}{r@{\;}l@{\;}l}
      \fun{rev}\,[3;2;1]
      & \smashedrightarrow{\epsilon} & \fun{rcat}([3;2;1],\el)\\
      & \smashedrightarrow{\eta}     & \fun{rcat}([2;1],[3])\\
      & \smashedrightarrow{\eta}     & \fun{rcat}([1],[2;3])\\
      & \smashedrightarrow{\eta}     & \fun{rcat}(\el,[1;2;3])\\
      & \smashedrightarrow{\zeta}    & [1;2;3].
    \end{array}
  \end{equation*}

\end{frame}

% ------------------------------------------------------------------------
%
\begin{frame}{Exact cost / Efficient reversal}

  To draw the cost of \fun{rev}, it is sufficient to notice that the
  first argument of \fun{rcat} strictly decreases at each rewrite, so
  all evaluation traces have the shape \(\epsilon\eta^n\zeta\), so
  \(\C{\fun{rev}}{n} = n + 2\). The cost is \textbf{linear}, so
  \fun{rev} must always be used instead of \fun{rev\(_0\)}.

  \bigskip

  The reason for the inefficiency of \fun{rev\(_0\)} can be seen in
  the fact that rule~\(\delta\) produces a series of calls to
  \fun{cat} following the pattern
  \begin{equation*}
    \fun{rev}_0\,s \twoheadrightarrow \fun{cat}(\fun{cat}(\dots
    \fun{cat}(\el, [x_n]), \dots, [x_2]), [x_1]),
  \end{equation*}
  where \(s = [x_1, x_2, \dots, x_n]\). The cost of all these calls to
  \fun{cat} is thus
  \begin{equation*}
    1 + 2 + \dots + (n-1) = \frac{n(n-1)}{2}.
  \end{equation*}
  because the cost of \(\fun{cat}(s,t)\) is \(n+1\). The problem is
  not calling \fun{cat}, but the fact that the calls are embedded in
  the most unfavourable configuration.

\end{frame}

% ------------------------------------------------------------------------
%
\begin{frame}{Extremal costs}

  When considering sorting programs based on comparisons, the cost
  varies depending on the algorithm and it also often depends on the
  original partial ordering of the keys.

  \bigskip

  Therefore, size does not capture all aspects needed to assess
  efficiency.

  \bigskip

  This quite naturally leads to consider bounds on the cost: for a
  given input size, we seek the configurations of the input that
  minimise and maximise the cost, respectively called the
  \begin{itemize}

    \item \textbf{minimum cost} (cost for the \textbf{best case}) and

    \item \textbf{maximum cost} (cost for the \textbf{worst case}).

  \end{itemize}

  \bigskip

  For example, some sorting algorithms have their worst case when the
  keys are already sorted, others when they are sorted in reverse
  order, etc.

\end{frame}

% ------------------------------------------------------------------------
%
\begin{frame}{Average cost}

  Once we obtain bounds on a cost, the question about the
  \textbf{average} or \textbf{mean cost} (also called \textbf{expected
    cost}) arises as well.

  \bigskip

  It is the arithmetic mean of the costs for all possible inputs of a
  given size.

  \bigskip

  Some care is necessary, as there must be a finite number of such
  inputs.

  \bigskip

  For instance, to assess the mean cost of sorting algorithms based on
  comparisons, it is usual to assume
  \begin{enumerate}

    \item that the input is a series of \(n\)~\textbf{distinct keys} and

    \item that the sum of the costs is taken over all their
      \textbf{permutations}, thus divided by~\(n!\), the number of
      permutations of size~\(n\).

  \end{enumerate}

\end{frame}

% ------------------------------------------------------------------------
%
\begin{frame}{Average cost}

  The uniqueness constraint actually allows the analysis to
  equivalently, and more simply, consider the permutations of
  \((1,2,\dots,n)\).

  \bigskip

  Some sorting algorithms, like \textbf{merge sort} or
  \textbf{insertion sort}, have their average cost
  \textbf{asymptotically equivalent} to their maximum cost.

  \bigskip

  In other words, for increasingly large numbers of keys, the ratio of
  the two costs become arbitrarily close to~\(1\).

  \bigskip

  Some others, like \textbf{Hoare's sort}, also known as
  \textbf{quicksort}, have the growth rate of their average cost being
  of a lower magnitude than the maximum cost, on an asymptotic scale.

\end{frame}

% ------------------------------------------------------------------------
%
\begin{frame}{Amortised cost}

  \begin{itemize}

    \item Sometimes an update is costly because it is delayed by an
      imbalance in the data structure that calls for an immediate
      remediation, but this remediation itself may lead to a state
      such that subsequent operations are faster than if the costly
      update had not happen.

    \item Therefore, when considering a series of updates, it may be
      overly pessimistic to sum the maximum costs of all the
      operations considered in isolation.

    \item Instead, \textbf{amortised analysis} takes into account the
      interactions between updates, so a lower maximum bound on the
      cost is derived.

    \item Note that this kind of analysis is inherently different from
      the average case analysis, as its object is the composition of
      different functions instead of independent calls to the same
      function on different inputs. \textbf{Amortised analysis is a
        worst case analysis of a sequence of updates.}

  \end{itemize}

\end{frame}
